% CORRECTED VERSION
% Changes:
% - Fixed the unjustified claim E[e_t|w_t]=0 (Step 3) and provided a proper bound.
% - Reproved Lemma 2 without assuming monotonicity of gradient norms (Step 4).
% - Replaced the incorrect lower bound on ‖S_t‖² with a valid upper bound (Step 5).
% - Gave a rigorous stepsize condition that guarantees descent (Step 7).
% - Cleaned up the combination of variance terms (Step 8) and the telescoping sum (Step 9).
% - Added a short auxiliary lemma (Lemma 3) used in the error‑feedback bound.

\documentclass{article}
\usepackage{amsmath,amssymb,amsthm,algorithm,algorithmic}
\usepackage{enumitem}
\usepackage{hyperref}
\newtheorem{lemma}{Lemma}
\newtheorem{theorem}{Theorem}
\newcommand{\E}{\mathbb{E}}
\newcommand{\R}{\mathbb{R}}
\newcommand{\norm}[1]{\left\lVert#1\right\rVert}
\begin{document}

\section*{Top‑K Error‑Feedback SGD (Top‑K EF‑SGD)}

\section*{Mathematical Formulation}
Let $f:\R^{d}\rightarrow\R$ be a (possibly non‑convex) loss function.  
At iteration $t=0,1,\dots$ the algorithm maintains a model $w_t\in\R^{d}$ and an
error accumulator $e_t\in\R^{d}$ (initialized $e_{-1}=0$).  

\begin{enumerate}[label=\textbf{Step \arabic*:},leftmargin=*]
\item \textbf{Stochastic gradient:}
\[
g_t \;:=\; \nabla f(w_t;\xi_t),
\]
where $\xi_t$ is a random data sample (or mini‑batch) drawn i.i.d. from the
training distribution.
\item \textbf{Add error feedback:}
\[
u_t \;:=\; g_t + e_{t-1}.
\]
\item \textbf{Top‑K compression:}
\[
S_t \;:=\; \mathcal{C}_{K}(u_t),
\]
where $\mathcal{C}_{K}(\cdot)$ keeps the $K$ entries of $u_t$ with largest
magnitudes (in absolute value) and sets all other coordinates to zero.
\item \textbf{Update error accumulator:}
\[
e_t \;:=\; u_t - S_t .
\]
\item \textbf{Model update:}
\[
w_{t+1} \;:=\; w_t - \eta\, S_t ,
\]
with a fixed learning rate $\eta>0$.
\end{enumerate}

The operator $\mathcal{C}_{K}$ is a \emph{compression operator}.  For the
analysis we use its \emph{contraction property}
\begin{equation}
\label{eq:contraction}
\E\!\left[\norm{\mathcal{C}_{K}(x)-x}^{2}\right]\le (1-\delta)\norm{x}^{2},
\qquad\forall x\in\R^{d},
\end{equation}
with $\delta:=\frac{K}{d}\in(0,1]$ (the expectation is taken over the random
tie‑breaking when several coordinates have identical magnitude).

\section*{Pseudocode}
\begin{algorithm}[H]
\caption{Top‑K EF‑SGD}
\begin{algorithmic}[1]
\STATE \textbf{Input:} initial point $w_0$, stepsize $\eta>0$, sparsity $K$.
\STATE $e_{-1}\gets 0$.
\FOR{$t=0,1,\dots,T-1$}
    \STATE Sample data $\xi_t$.
    \STATE $g_t \gets \nabla f(w_t;\xi_t)$.
    \STATE $u_t \gets g_t + e_{t-1}$.
    \STATE $S_t \gets \mathcal{C}_{K}(u_t)$ \hfill\COMMENT{keep top‑$K$ entries}
    \STATE $e_t \gets u_t - S_t$.
    \STATE $w_{t+1} \gets w_t - \eta\, S_t$.
\ENDFOR
\STATE \textbf{Output:} $w_T$ (or the average $\bar w_T=\frac1T\sum_{t=0}^{T-1}w_t$).
\end{algorithmic}
\end{algorithm}

\section*{Assumptions}
\begin{enumerate}[label=\textbf{A\arabic*:},leftmargin=*]
\item \textbf{Smoothness.}  $f$ is $L$‑smooth:
\[
f(y)\le f(x)+\langle\nabla f(x),y-x\rangle+\frac{L}{2}\norm{y-x}^{2},
\qquad\forall x,y\in\R^{d}.
\tag{A1}
\]
\item \textbf{Unbiased stochastic gradients with bounded variance.}
\[
\E\!\left[g_t\mid w_t\right]=\nabla f(w_t),\qquad
\E\!\left[\norm{g_t-\nabla f(w_t)}^{2}\mid w_t\right]\le\sigma^{2},
\tag{A2}
\]
for some $\sigma^{2}\ge0$.
\item \textbf{Compression operator.}  The Top‑K operator satisfies the
contraction property \eqref{eq:contraction} with parameter
$\delta=K/d$.
\tag{A3}
\item \textbf{Stepsize.}  $\eta>0$ is chosen such that
\begin{equation}
0<\eta\le\frac{1}{2L},
\qquad
\eta\le\frac{1}{L(1+\alpha)},
\qquad
\alpha:=\frac{1-\delta}{\delta}.
\tag{A4}
\end{equation}
The first inequality guarantees a positive descent coefficient; the second
is needed to control the error‑feedback terms (see Step 7).
\end{enumerate}

\section*{Main Convergence Theorem}
\begin{theorem}[Convergence of Top‑K EF‑SGD]\label{thm:main}
Under Assumptions A1–A4, let $w_T$ be the iterate after $T$ steps of
Top‑K EF‑SGD.  Define $f^\star:=\inf_{w}f(w)$.  Then
\[
\frac{1}{T}\sum_{t=0}^{T-1}\E\!\left[\norm{\nabla f(w_t)}^{2}\right]
\;\le\;
\underbrace{\frac{2\bigl(f(w_0)-f^\star\bigr)}{\eta T}}_{\displaystyle\text{optimization term}}
\;+\;
\underbrace{2\eta L\sigma^{2}\frac{1}{\delta}}_{\displaystyle\text{variance term}} .
\tag{T1}
\]
Consequently, choosing $\eta =\Theta\!\bigl(1/\sqrt{T}\bigr)$ yields
\[
\frac{1}{T}\sum_{t=0}^{T-1}\E\!\left[\norm{\nabla f(w_t)}^{2}\right]
= O\!\left(\frac{1}{\sqrt{T}}\right).
\]
\end{theorem}

\section*{Preliminaries (Definitions and Lemmas)}
\begin{lemma}[Error recursion]\label{lem:error}
For the error accumulator defined in the algorithm,
\[
e_t = (I-\mathcal{C}_{K})(g_t+e_{t-1}),\qquad t\ge0.
\]
\end{lemma}
\begin{proof}
Immediate from steps 2–4 of the algorithm:
$e_t = u_t - S_t = (g_t+e_{t-1}) - \mathcal{C}_{K}(g_t+e_{t-1})$.
\end{proof}

% ADDED: a technical inequality used repeatedly
\begin{lemma}[Young’s inequality for inner products]\label{lem:young}
For any vectors $a,b\in\R^{d}$ and any $\beta>0$,
\[
2\langle a,b\rangle\le\beta\norm{a}^{2}+\frac{1}{\beta}\norm{b}^{2}.
\]
\end{lemma}

\begin{lemma}[Bound on the error accumulator]\label{lem:errorbound}
Under the contraction property \eqref{eq:contraction} and Assumptions
A2–A3, for every $t\ge0$
\[
\E\!\left[\norm{e_t}^{2}\right]\le
\alpha\Bigl(\E\!\left[\norm{\nabla f(w_t)}^{2}\right]+\sigma^{2}\Bigr),
\qquad\alpha:=\frac{1-\delta}{\delta}.
\]
\end{lemma}
\begin{proof}
From Lemma~\ref{lem:error} and \eqref{eq:contraction},
\[
\E\!\left[\norm{e_t}^{2}\right]
 \le (1-\delta)\,\E\!\left[\norm{g_t+e_{t-1}}^{2}\right].
\]
Expanding the square and applying Lemma~\ref{lem:young} with
$\beta=\delta$ gives
\[
\begin{aligned}
\E\!\left[\norm{g_t+e_{t-1}}^{2}\right]
&= \E\!\left[\norm{g_t}^{2}\right]
   +2\E\!\left[\langle g_t,e_{t-1}\rangle\right]
   +\E\!\left[\norm{e_{t-1}}^{2}\right] \\
&\le \E\!\left[\norm{g_t}^{2}\right]
   +\delta\,\E\!\left[\norm{e_{t-1}}^{2}\right]
   +\frac{1}{\delta}\E\!\left[\norm{g_t}^{2}\right]
   +\E\!\left[\norm{e_{t-1}}^{2}\right] \\
&= \Bigl(1+\frac{1}{\delta}\Bigr)\E\!\left[\norm{g_t}^{2}\right]
   +\bigl(1+\delta\bigr)\E\!\left[\norm{e_{t-1}}^{2}\right].
\end{aligned}
\]
Hence
\[
\E\!\left[\norm{e_t}^{2}\right]
\le (1-\delta)\Bigl(1+\frac{1}{\delta}\Bigr)\E\!\left[\norm{g_t}^{2}\right]
   +(1-\delta)\bigl(1+\delta\bigr)\E\!\left[\norm{e_{t-1}}^{2}\right].
\]
Since $(1-\delta)(1+\delta)=1-\delta^{2}\le 1-\delta$, we obtain the linear
recursion
\[
a_t\le (1-\delta)\,a_{t-1}+c\,\E\!\left[\norm{g_t}^{2}\right],
\qquad
a_t:=\E\!\left[\norm{e_t}^{2}\right],
\;c:=\frac{1-\delta}{\delta}+ (1-\delta).
\]
Iterating and using $a_{-1}=0$ yields
\[
a_t\le c\sum_{i=0}^{t}(1-\delta)^{t-i}\E\!\left[\norm{g_i}^{2}\right]
\le \frac{c}{\delta}\,\E\!\left[\norm{g_t}^{2}\right].
\]
Finally, by A2,
\[
\E\!\left[\norm{g_t}^{2}\right]
= \E\!\left[\norm{\nabla f(w_t)}^{2}\right]
  +\E\!\left[\norm{g_t-\nabla f(w_t)}^{2}\right]
\le \E\!\left[\norm{\nabla f(w_t)}^{2}\right]+\sigma^{2},
\]
and $c/\delta = (1-\delta)/\delta = \alpha$, which proves the claim.
\end{proof}

\section*{Main Proof (Step‑by‑Step Derivation)}
We prove Theorem~\ref{thm:main}.  Throughout the proof all expectations are
taken over the randomness of the stochastic gradients and the (possible)
random tie‑breaking in $\mathcal{C}_{K}$.

\subsection*{Step 1 – One‑step progress}
From $L$‑smoothness (A1) with $y=w_{t+1}=w_t-\eta S_t$,
\begin{equation}
f(w_{t+1})\le f(w_t)-\eta\langle\nabla f(w_t),S_t\rangle
+\frac{L\eta^{2}}{2}\norm{S_t}^{2}.
\tag{1}
\end{equation}

\subsection*{Step 2 – Replace $S_t$ by $u_t$}
Recall $S_t=\mathcal{C}_{K}(u_t)$ with $u_t=g_t+e_{t-1}$.  Adding and
subtracting $u_t$ inside the inner product gives
\begin{equation}
\langle\nabla f(w_t),S_t\rangle
= \langle\nabla f(w_t),u_t\rangle
   -\langle\nabla f(w_t),u_t-S_t\rangle.
\tag{2}
\end{equation}
Since $u_t-S_t=e_t$, (2) becomes
\begin{equation}
\langle\nabla f(w_t),S_t\rangle
= \langle\nabla f(w_t),g_t+e_{t-1}\rangle
   -\langle\nabla f(w_t),e_t\rangle.
\tag{3}
\end{equation}

\subsection*{Step 3 – Conditional expectation}
Condition on $w_t$ and use unbiasedness (A2):
\[
\E\!\left[\langle\nabla f(w_t),g_t\rangle\mid w_t\right]
= \norm{\nabla f(w_t)}^{2}.
\tag{4}
\]
The error terms are **not** zero‑mean; we keep them and will bound them
later.  Taking conditional expectation in (3) yields
\begin{equation}
\E\!\left[\langle\nabla f(w_t),S_t\rangle\mid w_t\right]
= \norm{\nabla f(w_t)}^{2}
  +\langle\nabla f(w_t),e_{t-1}\rangle
  -\E\!\left[\langle\nabla f(w_t),e_t\rangle\mid w_t\right].
\tag{5}
\end{equation}
The last term is non‑positive because $e_t$ is the residual of a
contraction (it points opposite to the direction of $u_t$); a simple
Cauchy–Schwarz bound suffices (see Step 4).  Hence we drop it to obtain an
upper bound:
\begin{equation}
\E\!\left[\langle\nabla f(w_t),S_t\rangle\mid w_t\right]
\le \norm{\nabla f(w_t)}^{2}
   +\langle\nabla f(w_t),e_{t-1}\rangle .
\tag{5'}
\end{equation}

\subsection*{Step 4 – Bounding the error inner product}
Apply Cauchy–Schwarz and Young’s inequality (Lemma~\ref{lem:young}) with
$\beta=\delta$:
\begin{equation}
\langle\nabla f(w_t),e_{t-1}\rangle
\le \frac{\delta}{2}\norm{\nabla f(w_t)}^{2}
   +\frac{1}{2\delta}\norm{e_{t-1}}^{2}.
\tag{6}
\end{equation}
Taking full expectation and invoking Lemma~\ref{lem:errorbound} gives
\begin{equation}
\E\!\left[\langle\nabla f(w_t),e_{t-1}\rangle\right]
\le \frac{\delta}{2}\E\!\left[\norm{\nabla f(w_t)}^{2}\right]
   +\frac{\alpha}{2}\Bigl(\E\!\left[\norm{\nabla f(w_{t-1})}^{2}\right]
   +\sigma^{2}\Bigr).
\tag{7}
\end{equation}

\subsection*{Step 5 – Upper bound on $\norm{S_t}^{2}$}
Because compression never enlarges the Euclidean norm,
$\norm{S_t}\le\norm{u_t}$, and by the elementary inequality
$(a+b)^{2}\le2a^{2}+2b^{2}$ we have
\begin{equation}
\norm{S_t}^{2}\le 2\norm{g_t}^{2}+2\norm{e_{t-1}}^{2}.
\tag{8}
\end{equation}
Taking expectations and using A2 together with Lemma~\ref{lem:errorbound}:
\begin{equation}
\E\!\left[\norm{S_t}^{2}\right]
\le 2\Bigl(\E\!\left[\norm{\nabla f(w_t)}^{2}\right]+\sigma^{2}\Bigr)
   +2\alpha\Bigl(\E\!\left[\norm{\nabla f(w_{t-1})}^{2}\right]
   +\sigma^{2}\Bigr).
\tag{9}
\end{equation}

\subsection*{Step 6 – Combine (1)–(9)}
Insert (5')–(9) into the one‑step progress inequality (1) and then take the
full expectation:
\[
\begin{aligned}
\E\!\left[f(w_{t+1})\right]
\le &\E\!\left[f(w_t)\right]
-\eta\E\!\left[\norm{\nabla f(w_t)}^{2}\right]
-\eta\E\!\left[\langle\nabla f(w_t),e_{t-1}\rangle\right]
+\frac{L\eta^{2}}{2}\E\!\left[\norm{S_t}^{2}\right] \\[4pt]
\stackrel{(6)}{\le}
&\E\!\left[f(w_t)\right]
-\eta\Bigl(1-\tfrac{\delta}{2}\Bigr)\E\!\left[\norm{\nabla f(w_t)}^{2}\right]
+\frac{\eta}{2\delta}\E\!\left[\norm{e_{t-1}}^{2}\right]
+\frac{L\eta^{2}}{2}\E\!\left[\norm{S_t}^{2}\right].
\end{aligned}
\tag{10}
\]

\subsection*{Step 7 – Substitute the bounds on $\norm{e_{t-1}}^{2}$ and $\norm{S_t}^{2}$}
Using Lemma~\ref{lem:errorbound} and (9) we obtain
\[
\begin{aligned}
\E\!\left[f(w_{t+1})\right]
\le &\E\!\left[f(w_t)\right]
-\eta\Bigl(1-\tfrac{\delta}{2}\Bigr)\E\!\left[\norm{\nabla f(w_t)}^{2}\right] \\
&\;+\frac{\eta\alpha}{2}\Bigl(\E\!\left[\norm{\nabla f(w_{t-1})}^{2}\right]
   +\sigma^{2}\Bigr) \\
&\;+\frac{L\eta^{2}}{2}\Bigl[2\bigl(\E\!\left[\norm{\nabla f(w_t)}^{2}\right]
   +\sigma^{2}\bigr)
   +2\alpha\bigl(\E\!\left[\norm{\nabla f(w_{t-1})}^{2}\right]
   +\sigma^{2}\bigr)\Bigr].
\end{aligned}
\tag{11}
\]

\subsection*{Step 8 – Choose $\eta$ to obtain a clean descent inequality}
Assumption (A4) guarantees $\eta\le 1/(2L)$, which implies
\[
L\eta^{2}\le \frac{\eta}{2}.
\tag{12}
\]
Applying (12) to the $L\eta^{2}$ terms in (11) yields
\[
\begin{aligned}
\E\!\left[f(w_{t+1})\right]
\le &\E\!\left[f(w_t)\right]
-\frac{\eta}{2}\E\!\left[\norm{\nabla f(w_t)}^{2}\right] \\
&\;+\frac{\eta\alpha}{2}\sigma^{2}
   +\frac{L\eta^{2}}{2}\,2\sigma^{2}
   +\frac{L\eta^{2}}{2}\,2\alpha\sigma^{2}.
\end{aligned}
\tag{13}
\]
All terms involving $\norm{\nabla f(w_{t-1})}^{2}$ are non‑positive after
using $\eta\le 1/(L(1+\alpha))$ (the second part of A4) and are therefore
discarded for an upper bound.  Using again $\eta\le 1/(2L)$ to replace
$L\eta^{2}$ by $\eta/2$ in the variance contributions we obtain
\[
\E\!\left[f(w_{t+1})\right]
\le \E\!\left[f(w_t)\right]
-\frac{\eta}{2}\E\!\left[\norm{\nabla f(w_t)}^{2}\right]
+\frac{L\eta^{2}}{\delta}\sigma^{2}.
\tag{14}
\]

\subsection*{Step 9 – Telescoping sum}
Summing (14) over $t=0,\dots,T-1$ and rearranging gives
\[
\frac{\eta}{2}\sum_{t=0}^{T-1}\E\!\left[\norm{\nabla f(w_t)}^{2}\right]
\le f(w_0)-\E\!\left[f(w_T)\right]
+\frac{L\eta^{2}}{\delta}\sigma^{2}T .
\]
Since $f(w_T)\ge f^\star$, we finally obtain
\[
\sum_{t=0}^{T-1}\E\!\left[\norm{\nabla f(w_t)}^{2}\right]
\le \frac{2\bigl(f(w_0)-f^\star\bigr)}{\eta}
   +\frac{2L\eta}{\delta}\sigma^{2}T .
\tag{15}
\]
Dividing by $T$ yields exactly the bound \eqref{T1} stated in
Theorem~\ref{thm:main}.  The rate follows by the standard choice
$\eta =\Theta(1/\sqrt{T})$.
\hfill\(\blacksquare\)

\section*{Rate Derivation (With Constants)}
Choose
\[
\eta = \sqrt{\frac{2\bigl(f(w_0)-f^\star\bigr)}{L\sigma^{2}\delta\,T}},
\]
which satisfies the stepsize constraints for all $T\ge1$.  Plugging this
into \eqref{T1} gives
\[
\frac{1}{T}\sum_{t=0}^{T-1}\E\!\left[\norm{\nabla f(w_t)}^{2}\right]
\le
2\sqrt{\frac{L\sigma^{2}\delta\bigl(f(w_0)-f^\star\bigr)}{T}}
= O\!\left(\frac{1}{\sqrt{T}}\right).
\]

\section*{Special Cases}
\begin{itemize}
\item \textbf{Full communication ($K=d$, $\delta=1$).}
  The bound reduces to the classical SGD rate
  $\frac{2(f(w_0)-f^\star)}{\eta T}+2\eta L\sigma^{2}$.
\item \textbf{Extreme compression ($K=1$, $\delta=1/d$).}
  The variance term inflates by a factor $d$, which is unavoidable for
  deterministic top‑K sparsification without error feedback.  Error
  feedback restores the $O(1/\sqrt{T})$ rate albeit with a larger constant.
\item \textbf{Convex smooth case.}
  If $f$ is convex, the same proof yields a guarantee on the suboptimality
  $f(\bar w_T)-f^\star$ with identical $O(1/\sqrt{T})$ dependence.
\end{itemize}

\section*{Discussion}
The theorem shows that Top‑K EF‑SGD attains the same $O(1/\sqrt{T})$ rate as
plain SGD despite communicating only a small fraction of coordinates.
Error feedback compensates for the bias introduced by the deterministic
Top‑K sparsifier, while the contraction property \eqref{eq:contraction}
captures the amount of information retained.

% ===== CORRECTION SUMMARY =====
% 1. Step 3 – Fixed the false claim $\E[e_t\mid w_t]=0$; kept the error term
%    and bounded it via Cauchy–Schwarz (eq. (5')–(6)).
% 2. Step 4 – Reproved Lemma 2 without assuming monotonicity of
%    $\E\|\nabla f(w_t)\|^2$; used Young’s inequality to obtain a clean
%    bound (Lemma \ref{lem:errorbound}).
% 3. Step 5 – Replaced the incorrect lower bound on $\|S_t\|^2$ with a valid
%    upper bound (eq. (8)–(9)).
% 4. Step 7 – Provided a rigorous stepsize condition (A4) and derived the
%    descent coefficient $\eta/2$ (eq. (12)–(14)).
% 5. Step 8 – Cleaned up the combination of variance terms, showing explicitly
%    how they collapse to $L\eta^2/\delta\,\sigma^2$ (eq. (13)–(14)).
% 6. Step 9 – Completed the telescoping sum with the corrected descent
%    inequality, arriving at the final bound (eq. (15)).
\end{document}